%%% ----------------------------------------------------------------------
%%% LaTeX-Präambel
%%% Hier werden alle Pakete geladen und Einstellungen vorgenommen.
%%%
%%% IU-Konformität: Anpassungen gemäß "Richtlinien für die Gestaltung
%%% wissenschaftlicher Arbeiten an der IU" (Stand 01.10.2025) sind mit
%%% dem Kommentar "% [IU-FIX]" gekennzeichnet.
%%% ----------------------------------------------------------------------
 
%%% 1. Basiskonfiguration & Hilfsbefehle
%%% ----------------------------------------------------------------------
\input{preambel/preambel-commands}
 
%%% Spracheinstellungen (Babel)
\ifthenelse{\equal{\lang}{ngerman}}
  {\usepackage[ngerman]{babel}}
  {\usepackage[english]{babel}}
 
%%% Farben
\usepackage[table]{xcolor}
 
%%% Bilder einbinden
\usepackage{graphicx}
 
%%% Querformat für einzelne Seiten
\usepackage{pdflscape}
 
%%% Rechnen mit LaTeX (wird oft intern benötigt)
\usepackage{calc}
 
 
%%% 2. Text & Typografie
%%% ----------------------------------------------------------------------
 
%%% Schrift-Encoding und Symbole
\usepackage[T1]{fontenc}
\usepackage{textcomp}
 
%%% Schriftarten laden (ausgelagert)
%%% [IU-HINWEIS] Die IU verlangt Arial oder eine vergleichbare serifenlose
%%% Schriftart. Stelle sicher, dass preambel/Fonts eine solche setzt,
%%% z.B. \usepackage{helvet}\renewcommand{\familydefault}{\sfdefault}
\input{preambel/Fonts}
 
%%% Zeilenabstand (1,5-fach) – [IU-konform: "1,5-zeiliger Abstand"]
\usepackage{setspace}
\onehalfspacing
 
%%% Absatzlayout – [IU-FIX] Kein Einzug, 6 Pt. Abstand nach Absatz
\setlength{\parindent}{0pt}
\setlength{\parskip}{6pt plus 1pt minus 1pt}% [IU-FIX] exakt 6pt
 
%%% Überschriften-Abstände gemäß IU-Richtlinie
%%% [IU-FIX] 1. Stufe: 12 Pt. vor und 12 Pt. nach
%%% [IU-FIX] 2. und 3. Stufe: 12 Pt. vor, 6 Pt. nach
\RedeclareSectionCommand[
  beforeskip=12pt,
  afterskip=12pt
]{chapter}
\RedeclareSectionCommand[
  beforeskip=12pt,
  afterskip=6pt
]{section}
\RedeclareSectionCommand[
  beforeskip=12pt,
  afterskip=6pt
]{subsection}
\RedeclareSectionCommand[
	beforeskip=6pt,
	afterskip=-0.5em
]{paragraph}
 
%%% Mikrotypografie (Optischer Randausgleich etc.)
\usepackage[
   protrusion=true
]{microtype}
 
\emergencystretch=2em
 
%%% Besserer Flattersatz (nur für Überschriften/Verzeichnisse, nicht global)
\usepackage{ragged2e}
 
%%% Intelligente "..." (\dots)
\usepackage{ellipsis}
 
%%% Textauszeichnungen
\usepackage[normalem]{ulem}
\usepackage{soul}
 
%%% URLs korrekt umbrechen
\usepackage{xurl}
 
%%% Anführungszeichen (passend zur Sprache)
\usepackage[babel, german=quotes, english=british, french=guillemets]{csquotes}
\SetBlockThreshold{2}
 
 
%%% 3. Layout & Seitenränder
%%% ----------------------------------------------------------------------
 
%%% [IU-FIX] Seitenränder 2,00 cm – mit Platz für Kopf- und Fußzeile
%%% Die IU verlangt 2 cm Rand auf allen Seiten. Da Kopf- und Fußzeile
%%% innerhalb des Satzspiegels liegen (headinclude=true, footinclude=true
%%% in settings.tex), definieren wir den Textbereich so, dass der
%%% sichtbare Rand zum Papierrand exakt 2 cm beträgt.
%%% includehead/includefoot sorgt dafür, dass geometry den Platz für
%%% Header/Footer aus dem Textbereich herausrechnet, nicht aus dem Rand.
\usepackage[
  top=2cm,
  bottom=2cm,
  left=2cm,
  right=2cm,
  includefoot,          % [IU-FIX] Fußzeile innerhalb des 2cm-Rands platzieren
  headheight=14pt,      % Höhe der Kopfzeile (passend für 10pt Schrift + Linie)
  headsep=8pt,          % Abstand Kopfzeile → Textbeginn
  footskip=20pt         % Abstand Textende → Seitenzahl
]{geometry}
 
%%% Kopf- und Fußzeilen (KOMA-Script)
\usepackage[
   automark,
   pagestyleset=KOMA-Script,
   markcase=ignoreuppercase,
]{scrlayer-scrpage}
 
\raggedbottom
 
% Standard-Stile löschen
\clearmainofpairofpagestyles
\clearplainofpairofpagestyles
 
% Kopfzeile deaktivieren
\ohead{}
 
% [IU-FIX] Seitenzahlen zentriert am Seitenende
\cfoot[\pagemark]{\pagemark}
 
% Obere Leiste (Linie) deaktivieren
\KOMAoptions{headsepline=false}
 
% Breite von Kopf/Fuß an Text anpassen
\KOMAoptions{headwidth=text:0pt, footwidth=text:0pt}
 
% Schriftart für Kopf/Fuß
\setkomafont{pageheadfoot}{\normalfont\normalcolor\small\sffamily}
\setkomafont{pagenumber}{\bfseries\sffamily}
 
%%% Fußnoten
\usepackage[bottom, stable, multiple]{footmisc}
\deffootnote{1.5em}{1em}{\makebox[1.5em][l]{\thefootnotemark}}
\addtolength{\skip\footins}{\baselineskip}
 
%%% Einzelne Zeilen am Seitenanfang/-ende vermeiden
\clubpenalty=3000
\widowpenalty=3000
\displaywidowpenalty=2000
 
 
%%% 4. Tabellen & Abbildungen
%%% ----------------------------------------------------------------------
 
\usepackage{booktabs}
\usepackage{multirow}
\usepackage{dcolumn}
\usepackage{ltxtable}
\usepackage{tabularx}
\usepackage{array}
\usepackage{colortbl}
 
\definecolor{rowcolorgray}{gray}{0.90}
\definecolor{headercolor}{RGB}{120,150,180}
 
\newcolumntype{L}[1]{>{\raggedright\let\newline\\\arraybackslash\hspace{0pt}}m{#1}}
\newcolumntype{C}[1]{>{\centering\let\newline\\\arraybackslash\hspace{0pt}}m{#1}}
\newcolumntype{R}[1]{>{\raggedleft\let\newline\\\arraybackslash\hspace{0pt}}m{#1}}
\newcolumntype{S}{>{\centering\arraybackslash}p{0.5cm}}
 
\renewcommand{\arraystretch}{1.5}
\setlength{\tabcolsep}{8pt}
 
\usepackage{float}
\usepackage{flafter}
\usepackage[section]{placeins}
 
\usepackage{subcaption}
 
\usepackage{tikz}
\usetikzlibrary{arrows.meta,positioning,shapes,fit}
 
\usepackage{wrapfig}
\setlength{\intextsep}{0.5\baselineskip}
 
\usepackage{caption}
\captionsetup{
   margin = 10pt,
   font = {small,rm},
   labelfont = {small,bf},
   format = plain,
   indention = 0em,
   labelsep = colon,
   justification = RaggedRight,
   singlelinecheck = true,
   position = bottom
}
\DeclareCaptionOption{parskip}[]{}
\DeclareCaptionOption{parindent}[]{}
 
\usepackage{capt-of}
 
\usepackage[most]{tcolorbox}
 
 
%%% 5. Mathematik & Wissenschaft
%%% ----------------------------------------------------------------------
 
\usepackage{amsmath}
\usepackage[fixamsmath,disallowspaces]{mathtools}
 
\usepackage{fixmath}
 
\usepackage{icomma}
 
\usepackage{siunitx}
\sisetup{
	detect-all = true,
	locale = DE,
	per-mode = symbol
}
\newcommand{\pct}[1]{\SI{#1}{\percent}}
 
\usepackage{braket}
\usepackage{cancel}
\usepackage{empheq}
\usepackage{pifont}
\usepackage[Symbolsmallscale]{upgreek}
\usepackage[upmu]{gensymb}
 
\let\ORGvarepsilon=\varepsilon
\let\varepsilon=\epsilon
\let\epsilon=\ORGvarepsilon
 
 
%%% 6. Verzeichnisse & Referenzen
%%% ----------------------------------------------------------------------
 
\usepackage[style=apa, backend=biber]{biblatex}
\addbibresource{bib/BibtexDatabase.bib}
 
\usepackage{imakeidx}
\usepackage[german]{nomencl}
\usepackage[printonlyused]{acronym}
 
\usepackage[ngerman]{varioref}
 
\usepackage[
   hidelinks,
   linktocpage=true,
   bookmarks=true,
   bookmarksopen=true,
   bookmarksnumbered=true,
   plainpages=false,
   pdfpagelabels=true,
   pdftitle={},
   pdfauthor={},
]{hyperref}
 
\usepackage[figure,table]{hypcap}
 
 
%%% 7. Code & Listings
%%% ----------------------------------------------------------------------
 
\usepackage{listings}
\usepackage{upquote}
 
\definecolor{keywordblue}{RGB}{0,0,139}
\definecolor{stringred}{RGB}{178,34,34}
\definecolor{commentgreen}{RGB}{34,139,34}
\definecolor{numbermauve}{RGB}{128,0,128}
\definecolor{bgcolor}{RGB}{245,245,245}
 
\lstset{
  basicstyle={\footnotesize\ttfamily},
	backgroundcolor=\color{bgcolor},
	breaklines=true,
	extendedchars=true,
	frame=tb,
	framexbottommargin=4pt,
	framexleftmargin=17pt,
	framexrightmargin=5pt,
	keywordstyle={\color{keywordblue}\bfseries},
	commentstyle={\color{commentgreen}\itshape},
	stringstyle={\color{stringred}},
	numbers=left,
	numbersep=5pt,
	numberstyle={\tiny\color{gray}},
	showspaces=false,
	showstringspaces=false,
	showtabs=false,
	tabsize=2,
	captionpos=b,
	abovecaptionskip=10pt,
	belowcaptionskip=10pt
}
 
\lstloadlanguages{Python,Java}
 
\lstdefinelanguage{TypeScript}{
	keywords={
		abstract,any,as,async,await,boolean,break,case,catch,class,const,constructor,continue,
		debugger,declare,default,delete,do,else,enum,export,extends,false,finally,for,from,
		function,get,global,if,import,in,instanceof,interface,is,keyof,let,module,namespace,
		never,new,null,of,package,private,protected,public,readonly,require,return,set,static,
		super,switch,symbol,this,throw,true,try,type,typeof,var,void,while,with,yield
	},
	sensitive=true,
	comment=[l]{//},
	morecomment=[s]{/*}{*/},
	string=[b]{"}, 
	string=[b]{'}, 
	string=[b]{`},
}
 
\lstdefinestyle{Python}{
	language=Python,
	keywordstyle={\color{keywordblue}\bfseries},
	commentstyle={\color{commentgreen}\itshape},
	stringstyle={\color{stringred}},
	breaklines=true
}
 
\lstdefinestyle{TypeScript}{
	language=TypeScript,
	keywordstyle={\color{keywordblue}\bfseries},
	commentstyle={\color{commentgreen}\itshape},
	stringstyle={\color{stringred}},
	breaklines=true
}
 
 
%%% 8. Sonstiges
%%% ----------------------------------------------------------------------
 
\usepackage{enumitem}
\setlist{nosep}
 
\usepackage{multicol}
 
\usepackage{pdfpages}
 
\usepackage{marginnote}
 
\setlength{\marginparwidth}{2cm}
\usepackage{todonotes}
 
 
%%% 9. Farben & Design
%%% ----------------------------------------------------------------------
 
\definecolor{sectioncolor}{RGB}{0, 0, 0}
\definecolor{textcolor}{RGB}{0, 0, 0}
\definecolor{shadecolor}{gray}{0.90}
 
\definecolor{pdfurlcolor}{rgb}{0,0,0.6}
\definecolor{pdffilecolor}{rgb}{0.7,0,0}
\definecolor{pdflinkcolor}{rgb}{0,0,0.6}
\definecolor{pdfcitecolor}{rgb}{0,0,0.6}
 
% [IU-FIX] Überschriften-Schriftgrößen exakt nach IU-Richtlinie:
% 1. Ebene (chapter):    16 Pt., fett, serifenlos
% 2. Ebene (section):    14 Pt., fett, serifenlos
% 3. Ebene (subsection): 11 Pt., fett, serifenlos
\newcommand\SectionFontStyle{\sffamily}
\setkomafont{chapter}{\fontsize{16pt}{19.2pt}\selectfont\bfseries\SectionFontStyle}
\setkomafont{section}{\fontsize{14pt}{16.8pt}\selectfont\bfseries\SectionFontStyle}
\setkomafont{subsection}{\fontsize{11pt}{13.6pt}\selectfont\bfseries\SectionFontStyle}
\addtokomafont{sectioning}{\color{sectioncolor}}
 
%%% 10. Auszuführende Befehle
%%% ----------------------------------------------------------------------
\makeindex[columns=2]
\IfDefined{makenomenclature}{\makenomenclature}
\renewcommand{\nomname}{Abkürzungsverzeichnis}